\documentclass[journal,10pt,twocolumn]{IEEEtran}
\usepackage{amsmath,amssymb}
\usepackage{booktabs}
\usepackage{graphicx}
\usepackage{cite}
\usepackage{url}
\usepackage{microtype}
\usepackage{array}
\usepackage{multirow}
\usepackage{xspace}
\usepackage[T1]{fontenc}

% ── Macros ───────────────────────────────────────────────────────────────────
\newcommand{\STER}{\textsc{ster}\xspace}
\newcommand{\Tstar}{$T^*$\xspace}
\newcommand{\deli}{$\delta_i$\xspace}

\begin{document}

% ── Title & Author ────────────────────────────────────────────────────────────
\title{Architectural Isolation as a Timing Safety Primitive for Edge AI
Medical Devices: Controlled Experimental Evidence on a Shared-Silicon Platform}

\author{Akul~Mallayya~Swami%
  \thanks{A.~M.~Swami is with Varian Medical Systems, A Siemens
  Healthineers Company, Palo Alto, CA, USA
  (e-mail: swami.akul@alumni.uml.edu).
  ORCID: 0009-0003-9549-5543.
  Portions of this manuscript were drafted with the assistance of Claude
  (Anthropic). The author reviewed, verified, and takes full responsibility
  for all content, experimental data, and claims.}}

\markboth{IEEE Embedded Systems Letters}{}
\maketitle

% ── Abstract ─────────────────────────────────────────────────────────────────
\begin{abstract}
Output correctness and timing reliability are independent safety properties
of edge AI inference pipelines, yet FDA pre-market validation does not
explicitly characterize either under deployment conditions.
This letter demonstrates their independence through a controlled
same-hardware experiment: identical MobileNetV2 models are evaluated
under identical adversarial load on two execution paths of the same
NVIDIA Jetson Orin Nano Super: a dedicated GPU accelerator
(TensorRT~FP16) and a general-purpose CPU (ONNX Runtime FP32).
Across five stressor classes and ${\approx}158{,}000$ inference
activations, both paths maintain
$\mathrm{STER}{=}0.0000$ (zero exceedances above $T^*{=}0.05$),
confirming output instability is not the dominant failure mode.
The paths diverge sharply on timing: the GPU path maintains latency
${<}15$\,ms under all conditions; the CPU path degrades $7.2{\times}$
under combined load ($14.5$\,ms to $104.0$\,ms mean,
$\mathrm{P99}{=}165.1$\,ms), breaching the 10\,Hz clinical cycle budget
by 65\%.
Because the same-hardware design eliminates chip-to-chip confounds,
this divergence is consistent with causal attribution to the dedicated
accelerator architecture.
The Safety-Threshold Exceedance Rate (STER), formally defined as the
proportion of inferences exceeding a clinically-motivated ${\pm}5\%$
output deviation tolerance, provides a mechanism to verify output
stability independently of timing, enabling isolation of timing failures
from functional correctness failures.
Joint STER\,+\,latency verification is proposed as a candidate
pre-market acceptance protocol under FDA Draft Guidance
FDA-2024-D-4488 (January 2025).
\end{abstract}

\begin{IEEEkeywords}
edge AI, inference stability, timing safety, architectural isolation,
Safety-Threshold Exceedance Rate (STER), Software as a Medical Device
(SaMD), pre-market validation, IEC~62304, ISO~14971, TensorRT,
ONNX Runtime.
\end{IEEEkeywords}

% ── I. Introduction ───────────────────────────────────────────────────────────
\section{Introduction}

FDA pre-market validation of AI-enabled medical devices evaluates model
accuracy (sensitivity, specificity, AUC) under controlled,
zero-contention conditions~\cite{fda_draft_guidance}.
This framework does not explicitly characterize whether inference outputs
remain within bounds, or whether inference timing is preserved within
clinical cycle budgets, when the device operates under realistic
deployment load.
A 2025 cross-sectional study of 950 FDA-cleared AI-enabled devices found
that diagnostic errors were the leading cause of AI device recalls, with
43\% occurring within the first year of clearance~\cite{lee2025recalls}.
IEC~62304 Clause~6.3.2 requires verification that software meets its
timing requirements under all foreseeable operating conditions~\cite{iec62304};
however, no standardized empirical method exists for characterizing
whether inference pipeline timing is preserved under realistic deployment
stress.

This letter addresses this gap through a controlled experiment.
Running identical models under identical stressors on both the GPU
accelerator and the general-purpose CPU of the same physical hardware
eliminates platform confounds and supports attribution of timing behavior
to the accelerator architecture.
This work makes four contributions:
\begin{enumerate}
  \item \textbf{Output/timing independence}: output correctness (STER\,=\,0.0000)
    is preserved on both paths under all conditions; timing reliability
    diverges sharply. Current accuracy-based validation may not detect
    this distinction ($\Delta\mathrm{acc}{<}1.0\%$ on both paths).
  \item \textbf{Controlled same-hardware evidence}: GPU path maintains
    STER\,=\,0.0000 and latency ${<}15$\,ms; CPU path maintains
    STER\,=\,0.0000 but degrades to 104\,ms mean / 165\,ms P99 under
    combined load, breaching the 10\,Hz clinical cycle budget by 65\%.
  \item \textbf{STER as output stability verification}: STER formally
    verifies that inference output distributions remain within a
    clinically-motivated tolerance band, independently of timing, enabling
    isolation of timing failures from correctness failures.
  \item \textbf{Joint pre-market protocol}: STER\,$\leq$\,STER$_{\max}$
    AND latency$_{\mathrm{P99}}$\,$\leq$\,$T_n$ under foreseeable stressor
    conditions is proposed as a reproducible acceptance criterion
    operationalizing FDA-2024-D-4488's robustness assessment requirement.
\end{enumerate}

% ── II. Related Work ──────────────────────────────────────────────────────────
\section{Related Work}

Prior edge inference work addresses throughput maximization~\cite{hao2023},
latency predictability under GPU partitioning~\cite{martin2026}, and
cross-device accuracy benchmarking~\cite{baller2021,tobiasz2023}.
These works measure latency \emph{or} output accuracy, not their
independence under shared-silicon contention.
Safety-oriented work on neural networks in medical devices~\cite{becker2019}
addresses model-level mechanisms but does not empirically characterize
resource-contention-induced timing failure as a distinct safety property
from output correctness.
No prior study has used a same-hardware controlled design to isolate the
architectural contribution of a dedicated accelerator to timing
safety~\cite{martin2026}.

% ── III. System Model ─────────────────────────────────────────────────────────
\section{System Model and Formal Definitions}

\subsection{Output Stability Under Identical Inputs}

Let $f$ denote a deployed inference model and $\mathbf{x}_i$ a fixed
test input ($i = 1,\ldots,N$).
Under a stressor condition $c$, let $\sigma(\mathbf{y}_i^c)$ denote the
softmax output and $\bar{\sigma}_i$ the per-image reference softmax
captured at zero load.
The per-inference output deviation is:
\begin{equation}
  \delta_i = \bigl\|\sigma(\mathbf{y}_i^c) - \bar{\sigma}_i\bigr\|_\infty
           = \max_k \bigl|\sigma(\mathbf{y}_i^c)_k - \bar{\sigma}_{i,k}\bigr|
  \label{eq:delta}
\end{equation}

\subsection{STER Definition}

\textbf{Definition~1} ($T^*$): $T^* = 0.05$, selected as conservative
relative to ISO~14971 Table~B.1 clinical diagnostic tolerance bands
(typically 10--20\%).

\textbf{Definition~2} (STER): For $N$ inference activations:
\begin{equation}
  \mathrm{STER} = \frac{1}{N}\sum_{i=1}^{N}
    \mathbf{1}\!\left[\delta_i > T^*\right]
  \label{eq:ster}
\end{equation}
STER verifies output distribution bounds, which accuracy
($\mathrm{argmax}(\sigma(\mathbf{y}))$, a single label) structurally
cannot confirm.
STER is proposed as a candidate pre-market metric, not an existing
FDA-specified standard.

\subsection{Robustness Confirmation Condition}

The joint pre-market acceptance condition is:
\begin{equation}
  \mathrm{STER}_c \leq \mathrm{STER}_{\max}
  \quad\text{AND}\quad
  \mathrm{latency}_{\mathrm{P99}} \leq T_n
  \label{eq:accept}
\end{equation}
where $\mathrm{STER}_{\max} = 0$ and $T_n = 100$\,ms (10\,Hz cycle budget).
STER and latency are orthogonal safety dimensions: a device may satisfy
(\ref{eq:accept}) fully (GPU path, E0--E5), satisfy only the STER term
(CPU path, E6), or fail both. Both must be verified.

\subsection{Architectural Model}

The Jetson Orin Nano Super uses shared LPDDR5: CPU, GPU, and memory
controller compete for the same bus.
The \emph{GPU execution path} (TensorRT FP16) offloads inference to the
1024-core Ampere GPU via dedicated DMA paths with QoS bandwidth
reservation, isolating inference from CPU scheduling pressure.
The \emph{CPU execution path} (ONNX Runtime FP32) runs on the
6-core Arm Cortex-A78AE cluster, directly competing with stressor
workloads for CPU time, cache, and DRAM bandwidth.
This architectural difference is hypothesized to produce divergent
timing behavior under contention while preserving output correctness on both paths; this is empirically tested in E1--E6.

% ── IV. Hardware & Protocol ───────────────────────────────────────────────────
\section{Hardware Platform and Experimental Protocol}

\subsection{Hardware Platform}

\begin{table}[!t]
\caption{Hardware Platform Summary}
\label{tab:hardware}
\centering
\small
\begin{tabular}{p{2.4cm}p{1.5cm}p{3.6cm}}
\toprule
\textbf{Component} & \textbf{Role} & \textbf{Key Specification} \\
\midrule
Jetson Orin Nano Super 8GB (GPU path)
  & Primary inference host
  & 1024-core Ampere GPU, TensorRT FP16, shared LPDDR5 \\
Jetson Orin Nano Super 8GB (CPU path)
  & Negative control (same hardware)
  & ONNX Runtime FP32, CPUExecutionProvider, Cortex-A78AE \\
nRF52840 DK + TI SensorTag
  & BLE network stressor
  & BLE 5.0, up to 6 simultaneous connections \\
GL-AX1800 WiFi~6 Router
  & RF environment
  & 2.4\,GHz co-channel interference \\
\bottomrule
\end{tabular}
\end{table}

Both execution paths run MobileNetV2~\cite{sandler2018} (224$\times$224
RGB, ImageNet 1000-class head) on the same 500-image fixed test set,
held constant across all experiments.
The GPU path runs TensorRT FP16~\cite{tensorrt} at a nominal 10\,Hz
($T_n = 100$\,ms); the CPU path runs MobileNetV2 FP32 via ONNX
Runtime~\cite{onnxruntime} on the same hardware.
Using the same physical device for both paths is the critical
methodological choice: it eliminates chip-to-chip variation and thermal
differences as confounds, supporting attribution of timing behavior to
the accelerator architecture.

\subsection{Experimental Protocol}

\begin{table}[!t]
\caption{Experiment Summary}
\label{tab:experiments}
\centering
\small
\begin{tabular}{p{0.5cm}p{1.8cm}p{4.6cm}p{0.55cm}}
\toprule
\textbf{ID} & \textbf{Stressor} & \textbf{Protocol} & \textbf{IEC} \\
\midrule
E0 & Baseline & Both paths, zero load. Per-image reference $\bar{\sigma}_i$ captured. 10 trials $\times$ 500 inferences. & B+C \\
\addlinespace
E1 & CPU stress & stress-ng at 25/50/75/100\% system-wide load. Both paths tested at each level. 10 trials $\times$ 500 inferences per level. & B \\
\addlinespace
E2--E4 & Isolation stressors & Memory pressure (25--90\% fill), GPU co-tenancy (0--4 workers), network I/O (BLE 0--6 conns\,+\,WiFi). GPU path only; STER\,=\,0.0000 across all conditions. 5--10 trials per level. & B/C \\
\addlinespace
E5 & Combined load & CPU 75\%\,+\,memory 50\%\,+\,BLE 4 conns\,+\,fio disk write. Both paths. 10 trials $\times$ 500 inferences. \emph{Worst-case deployment scenario.} & C \\
\addlinespace
E6 & CPU negative control & CPU path only (no GPU). Conditions matching E0\,+\,E1\,+\,E5. 10 trials $\times$ 500 inferences per condition. Primary contrastive result. & B \\
\bottomrule
\end{tabular}
\end{table}

For each trial, $N = 500$ inference activations are recorded.
For each activation $i$: (1)~fixed test input $\mathbf{x}_i$ is
presented; (2)~softmax $\sigma(\mathbf{y}_i)$ is recorded;
(3)~$\delta_i$ is computed per~(\ref{eq:delta}) against the per-image
zero-load reference; (4)~latency is recorded via \texttt{time.perf\_counter()}.
STER is computed per~(\ref{eq:ster}) over the $N$ activations.
Top-1 accuracy is computed as a control variable.

% ── V. Results ────────────────────────────────────────────────────────────────
\section{Results}

\subsection{E0: Baseline Calibration}

\begin{table}[!t]
\caption{E0 Baseline: Both Execution Paths}
\label{tab:e0}
\centering
\small
\begin{tabular}{lcccc}
\toprule
\textbf{Path} & \textbf{STER} & \boldmath$\delta_\text{mean}$ & \textbf{Lat.\,mean} & \textbf{Lat.\,P99} \\
\midrule
GPU (TensorRT FP16) & 0.0000 & 0.0181 & 14.8\,ms & $\approx$18\,ms \\
CPU (ONNX FP32)     & 0.0000 & 0.0000 & 14.5\,ms & 18.6\,ms \\
\bottomrule
\end{tabular}
\end{table}

Both paths establish matched zero-load baselines ($\Delta$latency
${<}0.3$\,ms), confirming controlled experimental conditions.
GPU floating-point variance produces $\delta_\text{mean}{=}0.0181$ but
zero exceedances above $T^*{=}0.05$; CPU ONNX Runtime is
arithmetically deterministic ($\delta{=}0.0000$ exactly).
Top-1 accuracy: 1.6\% on both paths (same model, same images,
same label mapping; accuracy values reflect ImageNet label remapping
and are not used as performance benchmarks).

\subsection{E1: CPU Stress and Timing Divergence}

\begin{table}[!t]
\caption{E1 CPU Stress: Both Paths (10 trials $\times$ 500 inferences per level)}
\label{tab:e1}
\centering
\small
\begin{tabular}{ccccc}
\toprule
\multirow{2}{*}{\textbf{CPU Load}}
  & \multicolumn{2}{c}{\textbf{GPU Path}}
  & \multicolumn{2}{c}{\textbf{CPU Path}} \\
\cmidrule(lr){2-3}\cmidrule(lr){4-5}
  & STER & Lat.\,mean & STER & Lat.\,mean \\
\midrule
25\%  & 0.0000 & ${\approx}$14.8\,ms & 0.0000 & 33.4\,ms \\
50\%  & 0.0000 & ${\approx}$14.8\,ms & 0.0000 & 57.8\,ms \\
75\%  & 0.0000 & ${\approx}$14.8\,ms & 0.0000 & 73.7\,ms \\
100\% & 0.0000 & ${\approx}$14.8\,ms & 0.0000 & 70.9\,ms \\
\bottomrule
\end{tabular}
\end{table}

Both paths maintain STER\,=\,0.0000 at all CPU load levels;
$\delta_\text{mean}$ is load-invariant on both paths.
Timing diverges immediately: the GPU path is invariant
(${\approx}14.8$\,ms regardless of load) while the CPU path
scales with contention, reaching $73.7$\,ms at 75\% load, already $5.1{\times}$ the zero-load baseline.
TensorRT FP16 executes entirely on the GPU; the 6-core Arm cores
saturated by stress-ng have no architectural path to disturb the
inference pipeline.
ONNX Runtime FP32 competes directly with the stressor for CPU time
slices and cache bandwidth, producing latency degradation without
affecting arithmetic output.

\subsection{E2--E4: Isolation Stressors (GPU Path)}

Memory pressure (25--90\% LPDDR5 fill), GPU co-tenancy (0--4
concurrent TensorRT workers), and network I/O (BLE 0--6 connections
with WiFi co-channel interference) all produce STER\,=\,0.0000 with
$\delta_\text{mean}$ indistinguishable from E0 baseline on the GPU path
($\delta_\text{mean}{=}0.0181\pm0.0002$ across all conditions, $n{>}25{,}000$
inferences per stressor class).
GPU DMA isolation, dedicated QoS bandwidth reservation, and independent
TensorRT execution contexts prevent all tested stressor classes from
perturbing inference outputs or latency.
A confirmed +88\% CPU interrupt rate increase at BLE conns\,=\,4
(389\,$\to$\,732 IRQ/min) produced no measurable effect on GPU path
timing.

\subsection{E5: Combined Realistic Load}

Under simultaneous CPU 75\% + memory 50\% fill + BLE 4 connections
+ fio disk write, the GPU path maintains STER\,=\,0.0000 and latency
${<}15$\,ms (mean $14.9$\,ms, $\delta_\text{mean}{=}0.0156\pm0.0001$)
across all 10 trials.
The CPU path maintains STER\,=\,0.0000 but reaches 104.0\,ms mean /
165.1\,ms P99, exceeding the 10\,Hz cycle budget by 65\%.

\subsection{E6: CPU Negative Control}

\begin{table}[!t]
\caption{E6 CPU-Only Negative Control: Jetson Orin Nano Super,
  ONNX Runtime FP32 (10 trials $\times$ 500 inferences per condition)}
\label{tab:e6}
\centering
\small
\begin{tabular}{lcccc}
\toprule
\textbf{Condition} & \textbf{STER} & \boldmath$\delta_\text{mean}$ & \textbf{Lat.\,mean} & \textbf{Lat.\,P99} \\
\midrule
Zero load    & 0.0000 & 0.0000 & 14.5\,ms & 18.6\,ms \\
CPU 25\%     & 0.0000 & 0.0000 & 33.4\,ms & 111.4\,ms \\
CPU 50\%     & 0.0000 & 0.0000 & 57.8\,ms & 115.8\,ms \\
CPU 75\%     & 0.0000 & 0.0000 & 73.7\,ms & 122.1\,ms \\
CPU 100\%    & 0.0000 & 0.0000 & 70.9\,ms & 117.0\,ms \\
\textbf{Combined} & \textbf{0.0000} & \textbf{0.0000} & \textbf{104.0\,ms} & \textbf{165.1\,ms}$^\dagger$ \\
\bottomrule
\multicolumn{5}{l}{$^\dagger$ Exceeds $T_n{=}100$\,ms cycle budget by 65\%.}
\end{tabular}
\end{table}

STER\,=\,0.0000 and $\delta{=}0.0000$ exactly across all 50,000
inferences: ONNX Runtime's CPUExecutionProvider uses deterministic
ARM NEON SIMD operations; scheduling pressure delays computation
without perturbing its arithmetic.
Latency degrades monotonically with stressor intensity, reaching
$7.2{\times}$ baseline under combined load.
P99\,=\,165.1\,ms under combined load means 1\% of inferences
exceed the clinical cycle budget by ${\geq}65\%$.
The GPU path, under identical combined load, maintains latency
${<}15$\,ms across all 10 trials.
E6 establishes that the dedicated GPU accelerator's contribution to
the safety profile is \emph{timing isolation}, not numerical stability; this is a property of deterministic floating-point execution regardless of accelerator.

\subsection{Summary}

\begin{table}[!t]
\caption{Summary: GPU vs.\ CPU Path Under Worst-Case Conditions}
\label{tab:summary}
\centering
\small
\begin{tabular}{lcccc}
\toprule
\textbf{Path / Condition} & \textbf{STER} & \boldmath$\delta_\text{mean}$ & \textbf{Lat.\,mean} & \textbf{Lat.\,P99} \\
\midrule
GPU E0 (baseline)     & 0.0000 & 0.0181 & 14.8\,ms & $\approx$18\,ms \\
GPU E5 (combined)     & 0.0000 & 0.0156 & 14.9\,ms & $\approx$18\,ms \\
CPU E0 (baseline)     & 0.0000 & 0.0000 & 14.5\,ms & 18.6\,ms \\
CPU E5/E6 (combined)  & 0.0000 & 0.0000 & 104.0\,ms & 165.1\,ms$^\dagger$ \\
\bottomrule
\multicolumn{5}{l}{$^\dagger$ Exceeds $T_n$ by 65\%. STER passes; deployment timing fails.}
\end{tabular}
\end{table}

% ── VI. Discussion ────────────────────────────────────────────────────────────
\section{Discussion}

\subsection{Controlled Design as Methodological Contribution}

The same-hardware design is the critical methodological choice.
Cross-platform comparisons (e.g., Jetson vs.\ a different board) cannot
establish causality because thermal, memory topology, and driver-stack
differences confound the result.
Running GPU and CPU paths on the same Jetson Orin Nano Super controls
for all hardware-level confounds; the timing divergence of $7.2{\times}$
under identical stressors is consistent with causal attribution to the
accelerator architecture.
Platform datasheets document architectural separation in principle;
this work provides empirical quantification across simultaneous
multi-class stressor combinations at clinical deployment intensities.

\subsection{Timing as the Dominant Failure Mode}

STER\,=\,0.0000 on both paths establishes that output instability is
not the failure mode under realistic resource contention on this
hardware.
The dominant failure mode is timing degradation, which is invisible
to STER, invisible to accuracy, and invisible to any output-only
evaluation method.
This result motivates explicit timing characterization as a
complementary pre-market verification property, not as a replacement
for output stability metrics.

\subsection{Pre-Market Validation Gap}

Standard accuracy-based FDA pre-market testing would pass both
execution paths without distinguishing their timing safety profiles:
$\Delta\mathrm{acc}{<}1.0\%$ under all conditions on both paths.
A manufacturer using accuracy-only testing would clear both a
GPU-accelerated and a CPU-only deployment as equivalent, but the CPU-only deployment breaches the 10\,Hz clinical cycle budget by 65\%
under combined load.
The joint acceptance condition~(\ref{eq:accept}) closes this gap
with reproducible, hardware-agnostic criteria.

\subsection{STER as Output Stability Certification}

STER certifies that the full output distribution $\sigma(\mathbf{y})$
remains within the $T^*$ tolerance band, a property that
$\mathrm{argmax}(\sigma(\mathbf{y}))$ (accuracy) structurally cannot
confirm.
A model whose top-1 prediction is unchanged but whose confidence
distribution shifts by 4.9\% passes accuracy testing but may
violate $T^*$; STER detects this; accuracy does not.
STER\,=\,0.0000 on both paths in this study confirms output stability
under all tested conditions and isolates timing as the sole
safety-relevant divergence.

\subsection{Limitations}

Results are specific to MobileNetV2 at 10\,Hz on Jetson Orin Nano Super
under JetPack~6.
Models with higher memory bandwidth utilization, non-deterministic
preprocessing, or non-linear latency-throughput profiles may exhibit
different behavior.
$T^*{=}0.05$ is conservative relative to ISO~14971 Table~B.1 but is
not clinically validated for any specific device class; the appropriate
$T^*$ is a manufacturer and clinical risk management decision.

% ── VII. Conclusion ───────────────────────────────────────────────────────────
\section{Conclusion}

This letter has demonstrated, through a controlled same-hardware
experiment, that output correctness and timing reliability are
independent safety properties of edge AI inference pipelines.
Running MobileNetV2 under identical adversarial stressors on the GPU
and CPU execution paths of the same Jetson Orin Nano Super, both
paths maintain STER\,=\,0.0000 across ${\approx}158{,}000$ inference
activations; the GPU path maintains latency ${<}15$\,ms while the CPU
path degrades $7.2{\times}$ to 104.0\,ms mean / 165.1\,ms P99 under
combined load, breaching the 10\,Hz clinical cycle budget.
Current accuracy-based FDA pre-market testing would pass both paths
without detecting this timing failure.
STER formally verifies output distribution bounds; latency verification
under contention is required for timing safety.
These results indicate that ensuring output correctness alone may be
insufficient for reliable clinical deployment; temporal behavior under
realistic system contention must be explicitly validated.
Joint STER\,+\,latency verification is proposed as a candidate
pre-market acceptance protocol, operationalizing FDA Draft Guidance
FDA-2024-D-4488 and consistent with IEC~62304 and ISO~14971.

% ── References ────────────────────────────────────────────────────────────────
\bibliographystyle{IEEEtran}

\begin{thebibliography}{15}

\bibitem{fda_draft_guidance}
FDA, ``Artificial Intelligence-Enabled Device Software Functions:
Lifecycle Management and Marketing Submission Recommendations,''
Draft Guidance, Docket FDA-2024-D-4488, Jan.\ 2025.

\bibitem{fda_gmlp}
FDA, ``Good Machine Learning Practice for Medical Device Development:
Guiding Principles,'' Oct.\ 2021.

\bibitem{hao2023}
J.~Hao, P.~Subedi, L.~Ramaswamy, and I.\,K. Kim,
``Reaching for the Sky: Maximizing Deep Learning Inference Throughput
on Edge Devices with AI Multi-Tenancy,''
\emph{ACM Trans.\ Internet Technol.}, vol.~23, no.~1, pp.~1--33,
Feb.\ 2023.

\bibitem{martin2026}
J.\,J.~Mart\'{\i}n, A.~Castillo, J.~Garc\'{\i}a, and F.\,J.~Cazorla,
``Performance Isolation for Inference Processes in Edge GPU Systems,''
arXiv:2601.07600, Jan.\ 2026.

\bibitem{baller2021}
S.\,P.~Baller, A.~Jindal, M.~Chadha, and M.~Gerndt,
``DeepEdgeBench: Benchmarking Deep Neural Networks on Edge Devices,''
in \emph{Proc.\ IEEE Int.\ Conf.\ Cloud Eng.\ (IC2E)}, 2021,
pp.~20--30.

\bibitem{tobiasz2023}
R.~Tobiasz \emph{et al.},
``Edge Devices Inference Performance Comparison,''
arXiv:2306.12093, Jun.\ 2023.

\bibitem{becker2019}
U.~Becker,
``Increasing Safety of Neural Networks in Medical Devices,''
in \emph{Proc.\ SAFECOMP Workshops}, LNCS vol.~11699,
Springer, 2019, pp.~91--101.

\bibitem{iec62304}
IEC~62304:2006+AMD1:2015, ``Medical device software: Software life cycle processes,'' IEC, Geneva, 2015.

\bibitem{iso14971}
ISO~14971:2019, ``Medical devices: Application of risk management to medical devices,'' ISO, Geneva, 2019.

\bibitem{stressng}
C.\,I.~King, ``stress-ng: A Tool to Load and Stress a Computer
System,'' GitHub, 2023. [Online]. Available:
https://github.com/ColinIanKing/stress-ng

\bibitem{holoscan}
NVIDIA Corporation, ``Holoscan SDK: A Platform for Medical Device
Developers,'' NVIDIA Developer Documentation, 2024. [Online].
Available: https://developer.nvidia.com/holoscan-sdk

\bibitem{sandler2018}
M.~Sandler, A.~Howard, M.~Zhu, A.~Zhmoginov, and L.-C.~Chen,
``MobileNetV2: Inverted Residuals and Linear Bottlenecks,''
in \emph{Proc.\ IEEE/CVF CVPR}, 2018, pp.~4510--4520.

\bibitem{tensorrt}
NVIDIA Corporation, ``TensorRT Developer Guide,'' NVIDIA Developer
Documentation, 2024. [Online]. Available:
https://docs.nvidia.com/deeplearning/tensorrt/developer-guide/

\bibitem{onnxruntime}
Microsoft Corporation, ``ONNX Runtime: Cross-Platform Inference
Accelerator,'' GitHub, 2024. [Online]. Available:
https://github.com/microsoft/onnxruntime

\bibitem{lee2025recalls}
B.~Lee \emph{et al.},
``Early Recalls and Clinical Validation Gaps in Artificial
Intelligence-Enabled Medical Devices,''
\emph{JAMA Health Forum}, vol.~6, no.~8, p.~e253172, Aug.\ 2025.

\end{thebibliography}

\end{document}
